\chapter{\xxx{Background}}\label{chap:background}

Before discussing specific approaches to code optimization, we first outline the hardware features of modern high-performance systems that are most relevant to our focus on layout and traversal optimizations.

We target the two main general-purpose architectures used in high-performance computing: multi-core CPUs and GPU accelerators.
% The main difference between them is that CPUs are primarily built for task processing with some support for data parallelism, while GPUs are optimized for massively data-parallel workloads that run thousands of threads concurrently~\cite{hennessy2019computer,cuda_programming_guide}.
CPUs are primarily designed for task processing, executing each active thread on an independent core, while GPUs are optimized for massively data-parallel workloads that run thousands of threads concurrently in SIMT (Single Instruction, Multiple Threads) fashion~\cite{hennessy2019computer,cuda_programming_guide}, where threads run in groups that execute the same instruction at the same time on different data. This distinction, however, is not strict, as modern CPU cores also support data parallelism through using specialized vector units that can, in a single instruction, operate on multiple data elements (SIMD)\xxx{; and CPU cores may run multiple threads concurrently (at a cost of higher resource consumption)}.

Both CPUs and GPUs have hierarchies of memory units with different latencies and bandwidths. In both architectures, a computation operates on a limited amount of register memory, which is the fastest, but also the most limited memory unit --- for example, it does not allow for indirect accesses to data (e.g., accessing elements of a matrix using indices stored in memory). On the opposite end of the hierarchy, both architectures have access to a large global memory that can be accessed by all threads, but with high latency and low bandwidth (compared to other memory units).

In between these two extremes, both CPUs and GPUs have multiple levels of cache memory that offer better latency and bandwidth depending on how close the caches are to the computation units the computation is running on --- generally, there is at least one cache level dedicated to each core, and one shared cache level for all cores. The main goal of caches is the same regardless of the architecture or cache level: to store data that is likely to be accessed in the near future, so that such accesses are served faster and, for the closest cache levels, with less interference with other threads.

These caches are designed on basis of the principle of locality, which states that programs tend to access data repeatedly in a short period of time (temporal locality) and access data stored in nearby memory locations (spatial locality)~\xxx{(CITE)}. Many algorithms either naturally exhibit good locality (e.g., traversing a row-major matrix by rows) or can be transformed to have better locality (e.g., by reordering computations in such a way that repeated accesses to the same data are closer in time).

In general, caching is transparent to users\xxx{, unless they use specific instructions that specify caching behavior (e.g., prefetching instructions)}. However, even without using such instructions, a difference between a na\"{\i}ve implementation and a more cache-friendly one can lead to significant improvements in execution time.

A simple example that will help us illustrate the impact of data layouts and traversal orders on performance due to caching as well as other hardware features is the matrix multiplication algorithm, which computes each element of the output matrix ($C$) as a dot product of a row from the first input matrix ($A$) and a column from the second input matrix ($B$). The algorithm can be implemented in \cpp as follows:

\begin{lstlisting}[language=C++, caption={Matrix multiplication}, label={lst:matrix-multiplication}]
float C[M][N]; // M rows, N columns
float A[M][K]; // M rows, K columns
float B[K][N]; // K rows, N columns

// iterating over output matrix C
for (int i = 0; i < M; ++i) {
    for (int j = 0; j < N; ++j) {
        // single dot product
        float sum = 0;
        for (int k = 0; k < K; ++k) {
            sum += A[i][k] * B[k][j];
        }
        C[i][j] = sum;
    }
}
\end{lstlisting}

In \cref{lst:matrix-multiplication}, three characteristic features can be immediately observed:
\begin{enumerate}
    \item Each element of the output matrix $C$ can be computed independently in a separate thread, making the algorithm highly parallelizable.
    \item The innermost loop accesses unique elements of the input matrices $A$ and $B$ in each iteration with no reuse of data across iterations, making the algorithm memory-inefficient due to poor temporal locality (and probably memory-bound --- meaning that its performance is limited by memory access latency and bandwidth rather than computation cost).
    \item The elements of input matrix $B$ are accessed in a non-contiguous way, making the algorithm memory-inefficient due to poor spatial locality.
\end{enumerate}

On both CPUs and GPUs, we might easily take advantage of the first feature by parallelizing the outer loops ($i$ and $j$ loops), assigning each parallel thread to compute a different (or a group of) element(s) of the output matrix $C$. This, on its own, can lead to a massive improvement in execution time, depending on the number of threads we can run in parallel, since there are no data dependencies between the computations of different elements of $C$.

Taking advantage of the second feature is slightly more complex. A common way to do this is to break the computation into partitions of the input and output matrices that can fit in the target cache level. If we partition the matrices into tiles of $T \times T$ elements (choosing $T$ such that the data of the tiles can fit in the target cache memory) and update the corresponding elements of the output matrix by computing the dot products of the corresponding rows and columns of the tiles, then each value in the input matrices can be reused up to $T$ times during the computation of each tile of the output matrix.

The third feature can be addressed by changing the data layout of the input matrix $B$ to column-major (transposing it). This way, the innermost loop will access contiguous memory locations in both input matrices, improving spatial locality and thus performance. In CPUs, this can also allow the use of vector units that operate on multiple contiguous data elements at the same time, and in GPUs, such a transformation can make loads from threads accessing subsequent elements of $B$ coalesced (allowing the GPU to serve them together in fewer memory transactions).

These three optimizations showcase parallelization, traversal order, and data layout optimizations, respectively. All of them can be applied to the same algorithm; however, their combination may be non-trivial, and the benefits of each optimization may have complex interactions with the benefits of the other optimizations.
In \cref{chap:optimization-methods}, we will discuss various tools and frameworks that help with applying these, and many other, optimizations automatically, or allow users to express the algorithms and data structures in such a way that makes it easier to apply these optimizations.


\section{Further GPU considerations}\label{sec:gpu-considerations}

GPU accelerators are among the most widely used hardware platforms in high-performance computing~\cite{top500}. Their importance has further increased with the rise of large language models and other deep learning workloads~\xxx{(CITE)}.

Since NVIDIA GPU architectures currently dominate the high-performance computing landscape~\cite{top500}, we will default to their nomenclature when describing the relevant hardware features and use them as a reference point.

The execution model of GPUs is based on a virtual hierarchy of threads organized into (possibly 3-dimensional) grids of thread blocks. This organization is important because a common practice is to manually replace each parallelized loop in the algorithm with a dimension of the thread grid --- this may be done after applying other optimizations, such as tiling (so that each thread block maps to a tile of the output matrix, for example).

The grid is specified by the programmer when launching a GPU kernel (a function that runs on the GPU) and then each thread block is automatically scheduled on available streaming multiprocessors (SMs) of the GPU\@ until all blocks are fully executed. Each SM then runs the threads of the block in consecutive groups of 32 (called warps) in lockstep (executing same instructions). Since all threads in a warp execute on the same computation unit, it is important to ensure that they do not diverge in their execution paths as well as to ensure that they access memory in regular patterns --- while GPUs can handle some divergence and irregular memory access patterns, they can lead to significant performance degradation.

This makes the choice of data layout and traversal order even more important on GPUs, as its execution model is tightly coupled with the traversal loops replaced by the thread grid dimensions.

Grouping threads into blocks also enables programmers to use shared memory, which is a user-managed cache shared by all threads in the same block (running on the same SM core), as well as fast synchronization primitives between the threads. For the threads in the same warps, gpus also provide warp-wide communication primitives that allow threads to exchange data even faster.
